\documentclass[10pt,a4]{examenEPN0}
\usepackage{comandosEPN}

\usepackage{empheq}

\newcommand{\itemcp}{\stepcounter{enumi}\item[ \textcolor{red}{\theenumi.$^*$}] }

\newcommand{\iitemcp}{\stepcounter{enumii}\item[ \textcolor{red}{\theenumii)$^*$}] }


\materia{Álgebra Lineal}
%\autor{Departamento de Formación Básica}
\autor{GR \hspace{1cm}.}
%\fecha{Semestre 2026-A}
\fecha{Nombre: \hspace{9cm} Cédula:}
\logouno[1.6cm]{Logo_EPN.eps}
\logodos[2.6cm]{Logo_DFB.eps}

\newcommand{\la}{\lambda}

%Datos
\newcommand{\datosAlgebra}[2]{
\EnExamen
{
\vspace{5mm}
    \noindent
    Nombres: \hrulefill\
    Carrera: \rule{4.5cm}{0.4pt}\
    \ifthenelse{\equal{#1}{}}%
    {Paralelo: \underline{\parbox{1.8cm}{\parbox{1cm}{}}}\\[3mm]}
    {Paralelo: \underline{\parbox{1.8cm}{\centering #1}}\\[3mm]}
    C\'edula: \rule{5cm}{0.4pt}\
    Firma: \hrulefill\
    Fecha: \underline{\parbox{2.1cm}{\centering #2}}
}}


% \newcommand{\indicacionesPrueba}{%
% \EnExamen{
% \begin{indicaciones}
%     \item La prueba tiene una duraci\'on  de \textbf{120 minutos.}
%     \item No se permite la utilizaci\'on de dispositivos electr\'onicos (celular, tablets, calculadora, etc.)
%     \item Leer detenidamente los enunciados, \textbf{justificar la respuesta para acceder al puntaje total}.
%     \item La prueba es completamente individual. Los estudiantes deberán dejar sus teléfonos celulares guardados y no podrán tenerlos consigo (por ejemplo, en los bolsillos de su ropa) durante toda la evaluación. En caso de detectarse algún acto de deshonestidad académica o el incumplimiento de esta norma, no podrán seguir rindiendo la prueba y su calificación será \textbf{cero}.
% \end{indicaciones}
% \pagestyle{empty}
% }}

\newcommand{\indicacionesPrueba}{%
\EnExamen{
\begin{indicaciones}
    \item La prueba tiene una duraci\'on  de \textbf{60 minutos.}
    \item Leer detenidamente los enunciados, \textbf{justificar la respuesta} para acceder al puntaje total.
    \item No se permite la utilizaci\'on de dispositivos electr\'onicos (celular, tablets, calculadora, etc.)
    \item La prueba es completamente individual
    \item No se permite portar teléfonos celulares, relojes inteligentes ni otros dispositivos electrónicos durante el examen, aun cuando permanezcan en los bolsillos. Estos dispositivos deberán guardarse en la mochila antes del inicio de la evaluación y esta deberá mantenerse cerrada durante todo el examen.
    \item He leído las indicaciones y acepto que el incumplimiento de estas se penaliza con \textbf{cero}. Firma:...................................

    
\end{indicaciones}
\pagestyle{empty}
}}

\newcommand{\indicacionesExamen}{
\EnExamen{\begin{indicaciones}
    \item La evaluación tiene una duraci\'on  de \textbf{120 minutos.}
    \item No se permite la utilizaci\'on de dispositivos electr\'onicos (celular, tablets, calculadora, etc.)
    %\item Leer detenidamente los enunciados, justificar la respuesta para acceder al puntaje total.
    \item El examen es completamente individual. En caso de detectarse alg\'un acto de deshonestidad acad\'emica no podr\'a seguir rindiendo la prueba y su calificaci\'on ser\'a \textbf{cero}.
\end{indicaciones}
}\pagestyle{empty}
}

%% Imprimir las una hoja A4, dos veces en A5
\newcommand{\imprimirDoble}
{%
\usepackage{pgfmorepages}
\pgfmorepagesloadextralayouts
\pgfpagesuselayout{repeated 2-up}[a4paper,landscape]
\thispagestyle{empty}
}


%Este sí se puede comentar
\newenvironment{solucion}{\color{blue}\proof[Solución]}{\endproof\color{black}}




% \renewcommand{\@autor}{
% ---\\
% \phantom{.} 
% ...}
% \makeatother

\examen{Prueba 02: solucionario}
\makeatletter
\renewcommand{\@autor}{Departamento de Formación Básica}
\renewcommand{\@fecha}{GR5E -- Jueves de 7 a 9 -- 05/11/2026}
\makeatother
\begin{document}
\encabezado
\noindent\textbf{Criterio general.} Se asigna crédito parcial por los pasos
correctos. Un error aritmético se penaliza donde se origina; los pasos
posteriores se valoran por su coherencia, sin repetir la penalización.
Se aceptan procedimientos equivalentes que respeten el método solicitado.

\section*{1. Operaciones por filas y determinante (2 puntos)}
\noindent\textbf{Rúbrica.}
\begin{itemize}[itemsep=2pt]
    \item Justificar el cambio de signo por el intercambio: 0,5 puntos.
    \item Aplicar el factor 3 por multiplicar una fila: 0,5 puntos.
    \item Justificar que sumar un múltiplo de otra fila no altera el determinante: 0,5 puntos.
    \item Obtener el valor y seleccionar la opción correcta: 0,5 puntos.
\end{itemize}
Partimos de $\det(A)=5$. El intercambio $F_1\leftrightarrow F_3$
cambia el signo y produce un determinante de $-5$. La operación
$F_2\leftarrow3F_2$ lo multiplica por 3, dando $-15$.
Finalmente, $F_4\leftarrow F_4+2F_1$ no cambia el determinante.
Por tanto, $\boxed{\det(M)=-15}$, opción \textbf{(b)}.

\section*{2. Precios de los servicios (6 puntos)}
\noindent\textbf{Rúbrica.}
\begin{itemize}[itemsep=2pt]
    \item Definir las variables y sus unidades: 0,5 puntos.
    \item Plantear las tres ecuaciones: 1,5 puntos (0,5 por ecuación).
    \item Identificar $A$, $x$ y $b$: 1 punto.
    \item Construir $[A\mid I_3]$: 0,25 puntos.
    \item Escalonar identificando las operaciones por filas: 1 punto.
    \item Reducir hasta $[I_3\mid A^{-1}]$: 1 punto.
    \item Calcular $A^{-1}b$: 0,5 puntos.
    \item Interpretar los precios: 0,25 puntos.
\end{itemize}
\noindent\textbf{a) Modelo (3 puntos).}
Sean $r$ el precio de reacondicionar un envase, $e$ el precio de una entrega
y $d$ el depósito por envase, todos en dólares. El gasto definitivo de la
primera clienta excluye el depósito devuelto. Así,
\[
    \begin{cases}r+e+d=95,\\r+e=65,\\2r+e=105.\end{cases}
    \qquad
    A=\begin{pmatrix}1&1&1\\1&1&0\\2&1&0\end{pmatrix},\quad
    x=\begin{pmatrix}r\\e\\d\end{pmatrix},\quad
    b=\begin{pmatrix}95\\65\\105\end{pmatrix}.
\]
\noindent\textbf{b) Inversa mediante Gauss--Jordan (3 puntos).}
Aplicamos cada operación a toda la fila de $[A\mid I_3]$.
Para escalonar, primero anulamos las entradas debajo del primer pivote:
\[
    \left[\begin{array}{ccc|ccc}
        1&1&1&1&0&0\\1&1&0&0&1&0\\2&1&0&0&0&1
    \end{array}\right]
    \xrightarrow{\substack{F_2\leftarrow F_2-F_1\\F_3\leftarrow F_3-2F_1}}
    \left[\begin{array}{ccc|ccc}
        1&1&1&1&0&0\\0&0&-1&-1&1&0\\0&-1&-2&-2&0&1
    \end{array}\right].
\]
Intercambiamos las filas para obtener la forma escalonada:
\[
    \xrightarrow{F_2\leftrightarrow F_3}
    \left[\begin{array}{ccc|ccc}
        1&1&1&1&0&0\\0&-1&-2&-2&0&1\\0&0&-1&-1&1&0
    \end{array}\right]
    \xrightarrow{\substack{F_2\leftarrow-F_2\\F_3\leftarrow-F_3}}
    \left[\begin{array}{ccc|ccc}
        1&1&1&1&0&0\\0&1&2&2&0&-1\\0&0&1&1&-1&0
    \end{array}\right].
\]
Para reducir, anulamos las entradas sobre el tercer pivote y luego sobre el segundo:
\[
    \xrightarrow{\substack{F_2\leftarrow F_2-2F_3\\F_1\leftarrow F_1-F_3}}
    \left[\begin{array}{ccc|ccc}
        1&1&0&0&1&0\\0&1&0&0&2&-1\\0&0&1&1&-1&0
    \end{array}\right]
    \xrightarrow{F_1\leftarrow F_1-F_2}
    \left[\begin{array}{ccc|ccc}
        1&0&0&0&-1&1\\0&1&0&0&2&-1\\0&0&1&1&-1&0
    \end{array}\right].
\]
El bloque izquierdo es $I_3$; por tanto,
\[
    A^{-1}=\begin{pmatrix}0&-1&1\\0&2&-1\\1&-1&0\end{pmatrix},\qquad
    x=A^{-1}b=\begin{pmatrix}-65+105\\130-105\\95-65\end{pmatrix}
    =\begin{pmatrix}40\\25\\30\end{pmatrix}.
\]
\textbf{Reacondicionamiento: \$40 por envase; entrega: \$25; depósito: \$30 por envase.}
Se comprueba que $40+25+30=95$, $40+25=65$ y $2(40)+25=105$.

\newpage
\section*{3. Sistema de cuatro incógnitas y Cramer (6 puntos)}
\noindent\textbf{Rúbrica.}
\begin{itemize}[itemsep=2pt]
    \item Escribir la matriz $A$ en el orden $(x,y,z,w)$: 1 punto.
    \item Calcular $\det(A)$ con procedimiento: 1,5 puntos.
    \item Justificar la unicidad por $\det(A)\ne0$: 0,5 puntos.
    \item Construir $A^{(1)}$ y $A^{(3)}$: 1 punto (0,5 por matriz).
    \item Calcular sus determinantes: 1 punto (0,5 por determinante).
    \item Aplicar los cocientes de Cramer para $x$ y $z$: 1 punto (0,5 por incógnita).
\end{itemize}
\noindent\textbf{a) Matriz y determinante (3 puntos).}
En el orden $(x,y,z,w)$,
\[
    A=\begin{pmatrix}2&0&1&3\\0&2&4&0\\0&0&1&3\\0&0&3&2\end{pmatrix},
    \qquad b=\begin{pmatrix}-7\\2\\-7\\0\end{pmatrix}.
\]
Expandiendo sucesivamente por la primera columna,
\[
    \det(A)=2\det\begin{pmatrix}2&4&0\\0&1&3\\0&3&2\end{pmatrix}
    =4\det\begin{pmatrix}1&3\\3&2\end{pmatrix}
    =4(2-9)=-28\ne0.
\]
En consecuencia, el sistema tiene una única solución.

\noindent\textbf{b) Cramer para $x$ y $z$ (3 puntos).}
Sea $A^{(i)}$ la matriz que reemplaza la columna $i$ de $A$ por $b$:
\[
    A^{(1)}=\begin{pmatrix}-7&0&1&3\\2&2&4&0\\-7&0&1&3\\0&0&3&2\end{pmatrix},
    \qquad
    A^{(3)}=\begin{pmatrix}2&0&-7&3\\0&2&2&0\\0&0&-7&3\\0&0&0&2\end{pmatrix}.
\]
La primera y la tercera fila de $A^{(1)}$ son iguales, luego
$\det(A^{(1)})=0$. La matriz $A^{(3)}$ es triangular superior, de modo que
\[
    \det(A^{(3)})=2\cdot2\cdot(-7)\cdot2=-56.
\]
Por la regla de Cramer,
\[
    \boxed{x=\frac{\det(A^{(1)})}{\det(A)}=\frac{0}{-28}=0,
    \qquad z=\frac{\det(A^{(3)})}{\det(A)}=\frac{-56}{-28}=2.}
\]

\section*{4. Expansión por cofactores (6 puntos)}
\noindent\textbf{Rúbrica.}
\begin{itemize}[itemsep=2pt]
    \item Elegir una fila o columna y escribir la expansión con signos correctos: 2 puntos.
    \item Formar los menores de los términos no nulos: 2 puntos (1 por menor).
    \item Calcular los determinantes de orden 2: 1 punto (0,5 por determinante).
    \item Efectuar la suma y concluir: 1 punto.
\end{itemize}
Para $A=\begin{pmatrix}3&-1&-1\\2&1&0\\3&1&1\end{pmatrix}$,
expandimos por la tercera columna, cuyos signos son $+,-,+$:
\begin{align*}
    \det(A)&=(-1)(+1)\det\begin{pmatrix}2&1\\3&1\end{pmatrix}
        +0(-1)\det\begin{pmatrix}3&-1\\3&1\end{pmatrix}
        +1(+1)\det\begin{pmatrix}3&-1\\2&1\end{pmatrix}\\
    &=-1(2-3)+0+1(3+2)=1+5=\boxed{6}.
\end{align*}
\end{document}
